\section{Related Work}

\subsection{Synthetic Training Data for NLU}

A few papers investigate the use of synthetic data for training ASAS models. For instance, studies use LLMs to generate responses for underrepresented score levels~\citep{fangUsingGPT4Augment2024} or for new questions without annotated responses~\citep{bexteLunchFreeOvercoming2025}, finding that synthetic student answers are less realistic and are not helpful for training ASAS models. However, these studies rely mostly on prompting off-the-shelf LLMs and do not utilize the large number of real student responses to seen questions.

Beyond ASAS, data synthesis has been widely studied for other NLU tasks~\citep{fengSurveyDataAugmentation2021,longLLMsDrivenSyntheticData2024}.
Early work fine-tuned synthesizers on labeled examples with the class label prepended to the input and then sampled new examples for each class~\citep{anabytavorNotHaveEnough2020,kumarDataAugmentationUsing2020}.
With larger LLMs, later studies explored data synthesis by prompting with label-descriptive instructions and trained small task models on the generated data without any annotated examples~\citep{schickGeneratingDatasetsPretrained2021,yeZeroGenEfficientZeroshot2022,mengGeneratingTrainingData2022}.
Since synthetic data tends to be noisy and repetitive, some studies improve diversity through attribute-conditioned prompts, personas, or decoding interventions~\citep{yuLargeLanguageModel2023,chungIncreasingDiversityMaintaining2023,geScalingSyntheticData2025}.

Filtering the best synthetic data is another research focus. Methods range from simple length and formatting heuristics~\citep{yeZeroGenEfficientZeroshot2022} to LLM-based acceptance judgments~\citep{anikinaRigorousEvaluationLLM2025}. However, \citet{pecherBetterGeneratorsClassifiers2026} show that the gain from LLM-based filtering can still saturate as more synthetic data is added. Other approaches select examples by generation probability~\citep{mengGeneratingTrainingData2022} or using a classifier trained on real data~\citep{anabytavorNotHaveEnough2020}. \citet{zouFuseGenPLMFusion2024} measure how much several small task models disagree on the probability of a sample's synthetic label, and argue that a useful synthetic dataset contains both samples the models agree on, which aid convergence, and ambiguous samples on which they disagree, which improve generalization. They then use this signal to reweight the training data. Manual label correction can make additional synthetic answers useful for scoring~\citep{bexteLunchFreeOvercoming2025}. We train a panel of judges on the real training data and use their probability of the requested label as a quality signal to select synthetic answers for training.
% We investigate how judges trained on real student answers can guide selection for unseen questions and domains as the amount of synthetic data increases.

\citet{liSyntheticDataGeneration2023} further show that synthetic data is less effective for more subjective tasks, such as sarcasm detection. In ASAS, the boundaries between partially correct and fully correct or incorrect answers are often vague, so a successful synthesizer must capture partial understanding. 


\subsection{Pedagogical Role Simulation with LLMs} 
Though studies on synthesizing student responses with LLMs for downstream ASAS training are scarce, there is a growing body of work on simulating students or tutors with LLMs, predominantly in the context of mathematics tutoring dialogues.

For tutor simulation, \citet{macinaMathDialDialogueTutoring2023} paired human teachers with LLMs prompted to exhibit common student errors to collect the MathDial tutoring dataset.
For student simulation, researchers have conditioned simulated students on personality traits~\citep{liuPersonalityawareStudentSimulation2025} or cognitive levels~\citep{wuEmbracingImperfectionSimulating2025}.
Beyond dialogues, LLMs have been prompted to answer exam questions as students of different ability levels to estimate item difficulty~\citep{benedettoUsingLLMsSimulate2024} or to simulate student revisions for training feedback generators~\citep{nairClosingLoopLearning2024}.

Reinforcement learning has proven effective for both simulation tasks.
For tutors, simulated students provide reward signals, either via DPO on preference pairs ranked by a student model's predicted response correctness~\citep{scarlatosTrainingLLMbasedTutors2025} or via online training through multi-turn interactions with simulated students~\citep{dinucu-jianuProblemSolvingTeachingProblemSolving2025}.
As for student simulation, prompting alone cannot yield realistic answers: teachers who tutored the MathDial students reported that they \textbf{lacked authenticity}, used \textbf{overly complex} language, and were at times logically inconsistent~\citep{martynovaCanLLMsEffectively2025}.
Training on real student data helps: preference optimization yields far more realistic simulated students in tutoring dialogues than prompting~\citep{scarlatosSimulatedStudentsTutoring2026}, and DPO aligns simulated responses with student abilities estimated by item response theory (IRT) models~\citep{scarlatosSMARTSimulatedStudents2025}.
However, it is unclear whether such simulated student answers can be used to train ASAS models.
